\scanpage{293}
% archive sheet number: 294
% Vertical stroke to the left of the corollary.
\uline{Cor} 1

\noindent\begin{tabular}{|p{.91\linewidth}}
Les foncteurs $i_N,\tau_N$ commutent à \unsure{ttes|toutes} les sortes de limites inductives et projectives, $q_N$ commute aux $\varinjlim$, $\varphi_N$ aux $\varprojlim$.
Les foncteurs $i_N,\varphi_N$ sont pl. fid, les foncteurs $q_N,\deleted{\unsure{\varphi}}\tau_N$ sont des foncteurs localisation. \deleted{\unsure{Comme}} Pour $\tau_N$ (\unsure{puisque|comme} il est exact, et commute aux $\varinjlim$ et $\varprojlim$) c'est la localisation mod une sous-cat. abélienne épaisse, stable par $\varinjlim$ et $\varprojlim$ ; en fait, c'est précisément la sous-cat $\mathcal M^{N-1}$ (qui est ainsi décrite \deleted{\unsure{via}} en termes de $\mathcal M_N$).
\end{tabular}
